\documentclass[12pt]{article}
\usepackage[utf8]{inputenc}
\usepackage{listings}
\usepackage{xcolor}
\usepackage{geometry}
\usepackage{parskip}

\geometry{margin=2.5cm}

\definecolor{codebg}{RGB}{245,245,245}
\definecolor{keyword}{RGB}{0,112,0}
\definecolor{comment}{RGB}{128,128,128}
\definecolor{string}{RGB}{180,0,0}

\lstdefinestyle{pythonstyle}{
    backgroundcolor=\color{codebg},
    basicstyle=\ttfamily\small,
    keywordstyle=\color{keyword}\bfseries,
    stringstyle=\color{string},
    commentstyle=\color{comment}\itshape,
    language=Python,
    frame=single,
    rulecolor=\color{gray!40},
    breaklines=true,
    showstringspaces=false,
    tabsize=4,
    captionpos=b,
}

\title{\textbf{Wind Turbine Blade Fault Classification\\Using a Simple CNN}\\[0.5em]
\large Step 5 --- Preprocess and Augment Images}
\author{Amreen Batool}
\date{}

\begin{document}
\maketitle

\section*{Step 5 --- Preprocess and Augment Images}

\begin{lstlisting}[style=pythonstyle]
train_datagen = ImageDataGenerator(
    rescale=1./255,
    rotation_range=10,
    width_shift_range=0.1,
    height_shift_range=0.1,
    zoom_range=0.1,
    horizontal_flip=True
)

val_datagen  = ImageDataGenerator(rescale=1./255)
test_datagen = ImageDataGenerator(rescale=1./255)
\end{lstlisting}

\subsection*{Explanation}

\textbf{\texttt{rescale=1./255}}
\begin{itemize}
    \item Changes pixel values from the range 0--255 to 0--1.
\end{itemize}

\textbf{Training augmentation parameters:}
\begin{itemize}
    \item \texttt{rotation\_range=10} $\rightarrow$ rotates images slightly (up to $10^\circ$).
    \item \texttt{width\_shift\_range=0.1} $\rightarrow$ shifts image left or right by up to 10\%.
    \item \texttt{height\_shift\_range=0.1} $\rightarrow$ shifts image up or down by up to 10\%.
    \item \texttt{zoom\_range=0.1} $\rightarrow$ zooms in or out by up to 10\%.
    \item \texttt{horizontal\_flip=True} $\rightarrow$ randomly flips images horizontally.
\end{itemize}

\textbf{Note:} Validation and test data use \emph{only} rescaling — no augmentation — so that evaluation is performed on unmodified images.

\subsection*{Why Augmentation Is Used}

Data augmentation artificially increases the variety of training images by applying small random transformations. This helps:
\begin{itemize}
    \item Reduce overfitting (the model memorising training data).
    \item Improve generalisation to new, unseen images.
\end{itemize}

\end{document}
